\subsection{Traces}\label{sec:trace}
\subsubsection{Class I: Kinematic Consistency}\label{sec:class-i}

A trace $\mathcal{T}$ identifies 
a rule for associating one-dimensional fields to a given the continuum displacement \(\hat{\bm u}(x,\bm{r})\):
\begin{SaveEquation}{eq:trace-map}
\mathcal{T}: \hat{\bm{u}}(x, \bm{r}) \;\longmapsto\; 
\Big(\bm{u}(x),\, \bm{\theta}(x)
,\, \bm{\alpha}(x)
\Big)
% =:\bm{z}(x)
% \label{eq:trace-def}
\end{SaveEquation}
% such that the decomposition
% \begin{equation}
% \hat{\bm{u}}_{\rm P}(x, \bm{r}) = \bm{u}(x) + \bm{\theta}(x) \times \bm{r} 
% + \bm{\Phi}(\bm{r})\,\bm{\alpha}(x)
% \label{eq:trace-decomp}
% \end{equation}
% holds identically. 

There are infinitely many traces that may exactly represent a given solution \(\hat{\bm u}_{\rm P}\) to a Saint-Venant problem, and each imparts a distinct physical meaning to boundary conditions like \(\bm{\theta}_{\mathcal{T}}(x_0) = \mathbf{0}\) in a finite element simulation. 

% \begin{definition}[Axial Traces]
% \label{def:sv-trace}
% An axial \emph{trace} 
% on a continuum motion \(\hat{\bm u}\) 
% is a linear map
% \[
% \mathcal{T} :\CauchySpace\to C^1([0,L];\mathbb{R}^{6+n_w}),
% \qquad
% \hat{\bm{u}}(x,\bm{r}) \mapsto 
% \bm{z}(x) =
% \begin{pmatrix}
% \bm{u}(x),&
% \bm{\theta}(x),&
% \bm{\alpha}(x)
% \end{pmatrix},
% \]
% of the form
% $$[\mathcal{T}\hat{\bm{u}}](x) = \mathcal{T}_x\bigl[\hat{\bm{u}}(x, \cdot)\bigr]$$
% where $\{\mathcal{T}_x\}_{x \in I}$ is a family of linear functionals on section displacements, and \(\CauchySpace\) is the space of admissible motions \(\hat{\bm u}\). 
% A trace $\mathcal{T}$ is linear if, for any 
% $\hat{\bm{u}}_1,\hat{\bm{u}}_2 \in \CauchySpace$ and 
% $\alpha,\beta \in \mathbb{R}$,
% \[
% \mathcal{T}[\alpha \hat{\bm{u}}_1 + \beta \hat{\bm{u}}_2]
% =
% \alpha\,\mathcal{T}[\hat{\bm{u}}_1] + \beta\,\mathcal{T}[\hat{\bm{u}}_2].
% \]
% \end{definition}

% \begin{definition}[Trace-induced 1D strain map on $\mathcal U_{\rm SV}$]



% \begin{definition}[SV local parameter reconstruction from trace strains]
% \label{def:sv-local-reconstruction}
% Let $\mathcal T$ be Class I, with trace matrices $\TraceDeDpRoot\in\mathrm{GL}(6)$,
% $\TraceDpDeRoot=(\TraceDeDpRoot)^{-1}$, and $\TraceEvolveEE$ as in Prop.~\ref{prop:trace-matrix}.

% For an SV strain field $\bm e(x)$, define the reconstructed initial strain
% \[
% \bm e_0 \coloneq (\mathbf 1-x\TraceEvolveEE)\,\bm e(x)
% \qquad (\text{using }\TraceEvolveEE^2=\mathbf 0),
% \]
% and the reconstructed Saint~Venant parameter vector
% \begin{equation}
% \IesanVector_{\rm rec}(x)\coloneq \TraceDpDeRoot\,\bm e_0
% =\TraceDpDeRoot(\mathbf 1-x\TraceEvolveEE)\bm e(x).
% \label{eq:p-reconstruction}
% \end{equation}
% \end{definition}

% \subsection{Saint-Venant Motions}



It is natural to restrict attention to traces which exhibit the following properties:
\begin{enumerate}
    \item Strain measures \(\bm{e}\) generated under Saint-Venant loading are at most linear in the axial coordinate $\reflenx$. 
    \item Any Saint-Venant strain field is exactly reconstructed from only the plane section strains \(\bm{\gamma},\bm{\kappa}\).
    \item Continuum plane-section deformations are preserved, in the sense that section‑rigid fields are mapped to pure beam rigid motions.
\end{enumerate}
% It is useful to restrict attention to traces generating strains \(\bm{e}(x)\) which are at most linear in $x$ for any problem in the Saint-Venant class.
This ensures that the strains can be exactly represented by typical beam finite elements. 
This practical consideration is formalized as follows.

\begin{definition}[Class I]\label{def:class-1} % def:e1
A trace \(\mathcal{T}\) is called a \textbf{Class I} trace if:
\begin{enumerate}[label=(\roman*)]
    \item For all \(\hat{\bm u} \in \IesanMotions\oplus \RigidMotions\) the following relations are furnished by $x$-independent linear operators \(\TraceEvolveEE\),\(\TraceDeDpOne\), and \(\TraceStiffNM\in\mathrm{GL}(6)\):
    \begin{subequations}
    \begin{align}
        \bm{e}'(x) &=: \TraceEvolveEE\,\bm{e}(0) 
        \label{eq:req-t1} 
        \\ 
        \partial_{\IesanVector}\bm{e}'(x) &=: \TraceDeDpOne 
        \label{eq:req-d1}
        \\ 
        \partial_{\bm{e}} \snm(x) &=: \TraceStiffNM 
        \label{eq:req-c1}
    \end{align}
    \label{eq:class-i-a}
    \end{subequations}

    \item For all section-rigid deformations \(\hat{\bm u} =\bm{v}(x)+\bm{\omega}(x)\times\bm{r} \in \PlaneMotions\), the trace fields \(\bm{u},\bm{\theta} = \mathcal{T} [\hat{\bm u}]\) are:
    \begin{subequations}
    \begin{align}
        \bm{u}(x) &= \hat{\bm u}(x,\mathbf{0}) \\
        \bm{\theta}(x) &= \bm{\omega}(x)
    \end{align}
    \label{eq:req-p1}
    \end{subequations}
\end{enumerate}
\end{definition}
%
% These conditions 

Strains \(\bm{e}\) generated from a trace with the property \eqref{eq:class-i-a} can be written in the form:
\begin{equation}
\bm{e}(x) 
% = \TraceDeDpRoot(\mathbf{1} + x\TraceEvolveEP)\bm{p}
=(\mathbf{1} + x\TraceEvolveEE)\TraceDeDpRoot\bm{p}
\label{eq:def-T0}
\end{equation}

The equivalence class of traces with the properties \eqref{eq:class-i-a} is uniquely represented in the Saint~Venant class by a matrix \(\TraceDpDeRoot \in\mathrm{GL}(6)\), such that:
\begin{subequations}
    \begin{align} 
        \TraceEvolveEE &= \TraceDeDpRoot \IesanEvolveSP \TraceDpDeRoot \\
        \TraceDeDpOne &= \TraceDeDpRoot \IesanEvolveSP \label{eq:res-d1} \\
        \TraceStiffNM &= \IesanStiff \TraceDpDeRoot \label{eq:def-CT}
    \end{align}
\end{subequations}
where \(\IesanStiff\) is given by \eqref{eq:G0-neg-m} and 
\(\IesanEvolveSP\) is the evolution operator given by \eqref{eq:A-matrix-neg}.

% \begin{proposition}\label{prop:trace-matrix}
% The equivalence class of traces with the properties \eqref{eq:class-i-a} is uniquely represented in the Saint~Venant class by a matrix \(\TraceDpDeRoot \in\mathrm{GL}(6)\), and 
% \begin{subequations}
%     \begin{align} 
%         \TraceEvolveEE &= \TraceDeDpRoot \IesanEvolveSP \TraceDpDeRoot \\
%         \TraceDeDpOne &= \TraceDeDpRoot \IesanEvolveSP \label{eq:res-d1} \\
%         \TraceStiffNM &= \IesanStiff \TraceDpDeRoot \label{eq:def-CT}
%     \end{align}
% \end{subequations}
% where \(\IesanStiff\) is given by \eqref{eq:G0-neg-m} and 
% \(\IesanEvolveSP\) is the evolution operator given by \eqref{eq:A-matrix-neg}.
% \end{proposition}
% \begin{proposition}[Structure of Plane Traces]
% \label{prop:plane-trace-structure}
Furthermore, if $\mathcal T$ is Class I, then the trace matrix
$\TraceDeDpRoot$ must have the block form:
\begin{equation}
\TraceDeDpRoot = \begin{bmatrix}
1 & \mathbf{0} % \CenterTrace 
& 0 & \TraceNF^\top \\[4pt]
\mathbf{0} & \mathbf{0} & \TraceFT & \TraceFF \\[4pt] 
0 & \mathbf{0}^\top & 1 & \TraceTF^\top \\[4pt]
\mathbf{0} & -\FrameBasis^\times & \TraceBT & \TraceMF
\end{bmatrix}
\qquad\text{with}\qquad 
\operatorname{det}\TraceDeDpRoot = \operatorname{det}\left(\TraceFF - \TraceFT\otimes\TraceTF\right)
% = \left(\begin{array}{rrrr}
% & \FrameBasis\otimes(\FrameBasis+\CenterTrace) & & \FrameBasis\otimes\TraceNF + \TraceFT\otimes\FrameBasis + \TraceFF \\[4pt]
%  &  -\FrameBasis^\times & & \FrameBasis\otimes(\FrameBasis+ \TraceTF) + \TraceBT\otimes\FrameBasis + \TraceMF %\\[4pt]
% \end{array}\right)
\label{eq:T0-plane-structure}
\end{equation}
where \(\FrameBasis\cdot\TraceNF = \FrameBasis\cdot\TraceFT = \FrameBasis\cdot\TraceBT = \FrameBasis\cdot\TraceTF = 0\) and partitions correspond to 
$\bm{e} = (\epsilon, \bm{\gamma}_\perp, \varkappa, \bm{\kappa}_\perp)$ 
and $\bm{p} = (a_\varepsilon, \bm{a}_\Omega, a_\vartheta, \bm{b}_{\Section})$.
The free parameters comprise an 16-dimensional family\footnote{These parameters must satisfy the additional restriction that $\TraceDeDpRoot$ be invertible.}.
\iffalse
\begin{table}[H]
\begin{center}
\caption{Trace matrix block structure.}
\label{tab:T0-blocks}
\renewcommand{\arraystretch}{1.3}
\begin{tabular}{clll}
\toprule
Block & Size &  Units & Description \\ 
\midrule
$\TraceBT$ & $\mathbb{R}^2$  & $\sim$ 
& Curvature--torsion coupling \\
% $\CenterTrace$ & $\mathbb{R}^2$ &  L & Trace center location \\
$\TraceFT$ & $\mathbb{R}^2$  & L 
& Shear--torsion coupling \\
$\TraceTF$ & $\mathbb{R}^2$  & L %--- 
& Twist--flexure coupling \\
$\TraceMF$ & $\mathbb{R}^{2\times2}$  & \Length{1} % ---
& Curvature--flexure coupling \\
$\TraceNF$ & $\mathbb{R}^2$  & \Length{2} %L 
& Axial--flexure coupling \\
$\TraceFF$ & $\mathbb{R}^{2\times2}$  & \Length{2} %L   
& Shear--flexure coupling \\
\midrule
Total: & 16 & & \\
\bottomrule
\end{tabular}
\end{center}
\end{table}
\fi
% \end{proposition}

% \begin{proof}
This is demonstrated as follows. 
The plane section property \eqref{eq:req-p1} constrains how $\TraceDeDpRoot$ acts on the section-rigid part of each Saint Venant mode:

\textit{Extension} ($a_\varepsilon$): The section-rigid displacement is 
$x\,a_\varepsilon\FrameBasis$, giving $\bm{\theta} = \mathbf{0}$ and 
$\bm{u} = x\,a_\varepsilon\FrameBasis$. The strains are 
$\epsilon = a_\varepsilon$ with all other components zero. 
This fixes column 1 of $\TraceDeDpRoot$ to $(1, \mathbf{0}, 0, \mathbf{0})^\top$.

\textit{Bending} ($\bm{a}_\Omega$): The section-rigid rotation is 
$\bm{\theta} = -x\,\FrameBasis \times \bm{a}_\Omega$, giving curvature 
$\bm{\kappa}_\perp = -\FrameBasis \times \bm{a}_\Omega$. 

The shear strain 
$\bm{\gamma}_\perp$ receives no section-rigid contribution. 
This fixes columns 2--3: entries $(1, 2\text{--}3) = \CenterTrace^\top$, 
$(2\text{--}3, 2\text{--}3) = \mathbf{0}$, $(4, 2\text{--}3) = 0$, and 
$(5\text{--}6, 2\text{--}3) = -\FrameBasis^\times$.

\textit{Torsion} ($a_\vartheta$): The section-rigid rotation is 
$\bm{\theta} = x\,a_\vartheta\FrameBasis$, giving twist $\varkappa = a_\vartheta$. 
This fixes column 4 entries $(1, 4) = 0$, $(4, 4) = 1$, $(5\text{--}6, 4) = \mathbf{0}$.
The shear--torsion coupling $\TraceFT$ and curvature--torsion coupling $\TraceBT$
arise from warping and remain free.

\textit{Flexure} ($\bm{b}_{\Section}$): The section-rigid part contributes to 
$\epsilon$ and $\bm{\kappa}_\perp$ through the $x$-dependent terms, 
but these enter only in $\TraceDeDpOne = \TraceDeDpRoot\IesanEvolveSP$. 
All entries in columns 5--6 arise from warping contributions and remain free 
parameters: $\TraceNF$ (axial--flexure), $\TraceFF$ (shear--flexure), 
$\TraceTF$ (twist--flexure), and $\TraceMF$ (curvature--flexure).
% \end{proof}

Using \eqref{eq:res-d1} and \eqref{eq:T0-plane-structure} furnishes the following universal form for \(\TraceDeDpOne\):
\begin{equation}
\TraceDeDpOne =
\TraceDeDpRoot\IesanEvolveSP =
% \TraceDeDpRoot\TraceEvolveEP =
\begin{bmatrix}
0 & \mathbf{0} & 0 & -\CentroidLocation^\top %+\CenterTrace^{\top}
\\
\mathbf{0} & \mathbf{0} & \mathbf{0} & \mathbf{0} \\
0 & \mathbf{0} & 0 & \Zero \\
\mathbf{0} & \mathbf{0} & \mathbf{0} & -\FrameBasis^\times
\end{bmatrix}
% \label{eq:T1}
\label{eq:res-d1-block}
\end{equation}

Inverting \eqref{eq:T0-plane-structure} furnishes for \(\TraceDpDeRoot\):
\begin{SaveEquation}{eq:T0-inv}
\TraceDpDeRoot = \left[\begin{array}{cc|cc}
1 & \ShiftN^{\top} 
& \ShiftScaleN & \mathbf{0} \\
\mathbf{0} & \ShiftA & \ShiftF %\bm{\varsigma}_{12}^{\mathrm{F}} 
& \FrameBasis^\times \\[3pt] \hline &&& \\[-8pt]
0 & \ShiftT^\top & \ShiftScaleT & \mathbf{0} \\%[4pt]
\mathbf{0} & \ShiftB & \ShiftK %-\ShiftB\TraceFT 
& \mathbf{0}
\end{array}\right]
% \label{eq:T0-inv}
\end{SaveEquation}
with entries given in terms of the entries in \eqref{eq:T0-plane-structure} by :%\cref{tab:T0-inv-blocks}.

\iftrue
\begin{subequations}
\begin{align}
\ShiftB &=\left(\TraceFF-\TraceFT\otimes\TraceTF\right)^{-1} \\
\ShiftA &=\FrameBasis^\times\bigl(\TraceBT\otimes\TraceTF - \TraceMF\bigr)\,\ShiftB \\
\ShiftK &=-\ShiftB\TraceFT \\
\ShiftT &=-\ShiftB\TraceTF \\
\ShiftN &=- \ShiftB\TraceNF \\
\ShiftF&=\FrameBasis\times\Bigl(
-\TraceBT+\bigl(\TraceMF - \TraceBT\otimes\TraceTF\bigr)\,\ShiftB\TraceFT
\Bigr)\\
\ShiftScaleN&=\TraceNF\cdot\ShiftB\TraceFT \\
\ShiftScaleT &=1 - \TraceTF\cdot\ShiftK
\end{align}
\label{eq:Tinv-from-T}
\end{subequations}
\fi 

\iffalse
\begin{table}[H]
\begin{center}
\caption{Inverse trace matrix blocks.}
\label{tab:T0-inv-blocks}
\renewcommand{\arraystretch}{1.3}
\begin{tabular}{rlllll}
\toprule
\multicolumn{2}{@{}l}{Block} & Size & Units \\%& Coupling & \\ 
\midrule
$\ShiftB$ & $\!=\left(\TraceFF-\TraceFT\otimes\TraceTF\right)^{-1}$ & $\mathbb{R}^{2 \times 2}$ &  \Length{-2} 
  % & \(\bm\gamma\stackrel{\CenterShearLocation}{\rightarrow}\Twist\)
  % & $\bm{\gamma}\rightarrow\bm\alpha_{\WarpShear}\stackrel{\CenterShearLocation}{\rightarrow}\alpha_{\WarpTwistIesan}$ 
  \\
$\ShiftA$ & $\!=\FrameBasis^\times\bigl(\TraceBT\otimes\TraceTF - \TraceMF\bigr)\,\ShiftB$ & $\mathbb{R}^{2 \times 2}$ &  \Length{-1} 
  % &  $\bm\gamma\rightarrow\bm\kappa$
  % & --- 
  \\
$\ShiftK$& $\!=-\ShiftB\TraceFT$ & $\mathbb{R}^2$ &  \Length{-1} 
  % & $\Twist\rightarrow\Twist$ 
  % & $\Twist\rightarrow\bm\alpha_{\WarpShear}$ 
  \\
$\ShiftT$ & $\!=-\ShiftB\TraceTF$ & $\mathbb{R}^2$ &  \Length{-1} 
  % & --- & $\epsilon\rightarrow\alpha_{\WarpTwistIesan}$ 
  \\
% $\ShiftNK$ & $\!=-\FrameBasis\times\CenterTrace$ &  $\mathbb{R}^2$ &  L 
%   & --- & --- \\
$\ShiftN$ & $\!=- \ShiftB\TraceNF$ & $\mathbb{R}^2$ &  $\sim$ 
  % & $\bm\gamma\rightarrow\epsilon$ 
  % & ---
  \\
$\ShiftF$ 
  & $=\FrameBasis\times\Bigl(
-\TraceBT+\bigl(\TraceMF - \TraceBT\otimes\TraceTF\bigr)\,\ShiftB\TraceFT
\Bigr)$ & $\mathbb{R}^2$ &  $\sim$ 
  % & $\Twist\rightarrow\bm{\kappa}$ 
  % & --- 
  \\
$\ShiftScaleN$ & $\!=\TraceNF\cdot\ShiftB\TraceFT$ %-\CenterTrace\cdot\ShiftF 
& $\mathbb{R}$ &  L
  % & $\Twist\rightarrow\epsilon$ 
  % & --- 
  \\
$\ShiftScaleT$ & $\!=1 - \TraceTF\cdot\ShiftK$ & $\mathbb{R}$ &  $\sim$ 
  % & $\Twist\rightarrow\Twist$ 
  % & $\Twist\rightarrow\alpha_{\WarpTwistIesan}$
  \\
\bottomrule
\end{tabular}
\end{center}
\end{table}
\fi


In view of \eqref{eq:T0-inv} with \(\IesanStiff\) given by \eqref{eq:G0-neg-m}, the section stiffness \(\TraceStiffNM=\IesanStiff\TraceDpDeRoot\) takes the general form:
\iffalse
\begin{equation}
\TraceStiffNM 
=
\left[\begin{array}{cc|cc}
\AxialRigidity 
& \left(\AxialRigidity\ShiftN 
+ \ShiftA^{\top} EA\CentroidLocation\right)^{\top} 
& \left(EA\ShiftScaleN + EA\CentroidLocation\cdot \ShiftF\right) 
& EA\CentroidLocation^\top\FrameBasis^\times
\\
\mathbf{0} &
\ShearRigidity\,\ShiftB &
-\,\ShearRigidity\,\ShiftB\TraceFT &
\mathbf{0}
\\ \hline &&& \\[-8pt] 
0 &
\TwistRigidity\,(\ShiftT)^\top
& \TwistRigidity\,\ShiftScaleT &
\mathbf{0}
\\ %[8pt]
-\FrameBasis\times\AxialCoupling 
& -(\FrameBasis\times\AxialCoupling)\otimes\ShiftN
  -\FrameBasis^\times E\IntRoR\,\ShiftA 
&
 -\FrameBasis\times\AxialCoupling\,\ShiftScaleN
 -\FrameBasis\times E\IntRoR\,\ShiftF
&
% EA(\FrameBasis\times\CentroidLocation)\otimes\ShiftNK 
- \FrameBasis^\times E\IntRoR\,\FrameBasis^\times
\end{array}\right].
\label{eq:C-C1}
\end{equation}
\fi 
\begin{equation}
\TraceStiffNM 
=
\left[\begin{array}{cc|cc}
\AxialRigidity
& \left(\AxialRigidity\ShiftN 
+ \ShiftA^{\top}\AxialCoupling\right)^{\top} 
& \left(\AxialRigidity\ShiftScaleN + \AxialCoupling\cdot \ShiftF\right) 
& \AxialCoupling^\top\FrameBasis^\times
\\[2pt]
\mathbf{0} &
\TraceShearRigidity &
\TraceShearRigidity\TraceShearOrigin &
\mathbf{0}
\\ \hline &&& \\[-7pt] 
0 &
\TwistRigidity\,(\ShiftT)^\top
& \TwistRigidity\,\ShiftScaleT &
\mathbf{0}
\\ %[8pt]
-\FrameBasis\times\AxialCoupling 
& -(\FrameBasis\times\AxialCoupling)\otimes\ShiftN
  -\FrameBasis^\times E\IntRoR\,\ShiftA 
&
 -\FrameBasis\times\AxialCoupling\,\ShiftScaleN
 -\FrameBasis\times E\IntRoR\,\ShiftF
&
% EA(\FrameBasis\times\CentroidLocation)\otimes\ShiftNK 
- \FrameBasis^\times E\IntRoR\,\FrameBasis^\times
\end{array}\right].
\label{eq:C-C1}
\end{equation}
where:
\begin{align}
\TraceShearRigidity &\coloneq \ShearRigidity \ShiftB \\
\TraceShearOrigin &\coloneq-\TraceFT \\
\TraceTwistOrigin &\coloneq-\TwistRigidity (\TraceShearRigidity)^{-1}\ShiftB\TraceTF
\end{align}
The (2,2) block of \eqref{eq:C-C1} reveals the generalized shear correction coefficients:
\begin{equation}
\bm{K} = (\PlaneShearRigidity)^{-1}\ShearRigidity\ShiftB
\qquad\text{with}\qquad 
\PlaneShearRigidity \coloneq \int G\dA
\label{eq:shear-correction}
\end{equation}
\begin{remark}\label{rem:bend-trace-indep}
The (4,4) block of \eqref{eq:C-C1} is independent of the trace and therefore unambiguous.
\end{remark}
\begin{remark}
\eqref{eq:C-C1} is not generally symmetric. 
This is addressed in \cref{sec:class-iii}.
\end{remark}

% \clearpage
\subsubsection{Summary.}
The Class I traces represent a family of beam theories parameterized by 16 constants (\cref{tab:trace-consts}) which are all capable of exactly representing Saint~Venant solutions through Type 1 beam finite elements.

\begin{table}[H]
\begin{center}
\caption{Parameters identifying a trace in the Saint~Venant class.}
\label{tab:trace-consts}
\renewcommand{\arraystretch}{1.3}
\begin{tabular}{clll}
\toprule
Block & Size &  Units & Description \\ 
\midrule
$\TraceShearRigidity$ & $\mathbb{R}^{2\times2}$  & \Length{2} %L   
& Trace-consistent shear rigidity. \\
$\ShiftA$ & $\mathbb{R}^{2\times2}$  & \Length{1} % ---
& Curvature--flexure coupling \\
$\TraceShearOrigin$ & $\mathbb{R}^2$  & L 
& Shear--torsion coupling \\
$\TraceTwistOrigin$ & $\mathbb{R}^2$  & L %--- 
& Twist--flexure coupling \\
$\ShiftF$ & $\mathbb{R}^2$  & $\sim$ 
& Curvature--torsion coupling \\
$\ShiftN$ & $\mathbb{R}^2$  & $\sim$
& Axial--flexure coupling \\
\bottomrule
\end{tabular}
\end{center}
\end{table}
In terms of the independent constants of \cref{tab:trace-consts} the remaining blocks of the trace matrix \eqref{eq:T0-inv} are:
\begin{align}
\ShiftB &=(\ShearRigidity)^{-1}\TraceShearRigidity \\
\ShiftK &=(\ShearRigidity)^{-1}\TraceShearRigidity\TraceShearOrigin \\
\ShiftT &=(\TwistRigidity)^{-1}\TraceShearRigidity\TraceTwistOrigin \\
\ShiftScaleN &=\TraceNF\cdot\ShiftB\TraceFT \\
\ShiftScaleT &=1 - \TraceTF\cdot\ShiftK \\
\end{align}

% In a Type 3 finite element simulation, the trace matrix \(\TraceDpDeRoot\) is incorporated into the section response through the matrices \(\RigidShift\) \eqref{eq:class-ii-Le} and \(\ShapeShift\) \eqref{eq:class-i-La}.

% \begin{table}[H]
% \caption{Summary of symbols (Class I).}
% \label{tab:trace-symbols}
% \centering
% \begin{tabular}{@{}lllll@{}}
%     \toprule
%     Symbol &  & Definition & Comprises & \\ 
%     \midrule
%     $\snm$ & & & $N,\ShearForce,T,\bm{M}$ \\
%     $\mathbb{J}$ & & \eqref{eq:J-matrix} & $\FrameBasis^{\times}$ & \\ 
%     $\IesanStiff$ &  & \eqref{eq:G0-neg-m} & $\IntRoR$, $\CentroidLocation$ & \cref{tab:iesan-constants} \\ 
%     $\bm{p}$ & & \eqref{eq:def-iesan-vector} 
%       & \(a_{\varepsilon},\, \bm{a}_{\Section},\, a_{\vartheta},\, \bm{b}_{\Section}\) & \cref{tab:iesan-constants} \\
%     \LightRule
%     $\TraceDpDeRoot$ &  & \eqref{eq:T0-inv}  %\cref{prop:trace-matrix} %
%       & $\ShiftB$,$\ShiftA$,$\ShiftF$,$\ShiftK$,$\ShiftN$,$\ShiftNK$ 
%       & \cref{tab:T0-inv-blocks} \\
%     $\TraceEvolveEE$ & & \eqref{eq:def-T0} & \\ 
%     $\TraceStiffNM$ & & \eqref{eq:def-CT} & & \cref{tab:iesan-constants},\cref{tab:T0-inv-blocks}
%     \\[3pt] 
%      % & \eqref{eq:def-iesan-energy} 
%      \\
%     \bottomrule
% \end{tabular}
% \end{table}
